% ১৪.৩ C Program-এর গঠন
\section{C Program-এর গঠন}

\begin{learningbox}
এই টপিক শেষে শিক্ষার্থী একটি basic C program-এর অংশ চিনতে ও লিখতে পারবে।
\end{learningbox}

\begin{verbatim}
#include <stdio.h>

int main()
{
    int a, b, sum;
    scanf("%d %d", &a, &b);
    sum = a + b;
    printf("Sum = %d", sum);
    return 0;
}
\end{verbatim}

\begin{center}
\begin{tabular}{|p{0.28\textwidth}|p{0.52\textwidth}|}
\hline
\textbf{অংশ} & \textbf{কাজ} \\
\hline
\texttt{\#include <stdio.h>} & standard input-output library যুক্ত করে \\
\hline
\texttt{int main()} & program execution শুরু হয় \\
\hline
Declaration & variable-এর type ও name ঘোষণা করে \\
\hline
\texttt{scanf}/\texttt{printf} & input নেওয়া ও output দেখানো \\
\hline
\texttt{return 0;} & program সফলভাবে শেষ বোঝায় \\
\hline
\end{tabular}
\end{center}

\begin{boardbox}{বোর্ড ফোকাস}
Code লেখার প্রশ্নে indentation, semicolon, brace এবং format specifier ঠিক রাখা জরুরি।
\end{boardbox}
